\section{Related work, discussion and open problems}\label{sec:related}

\paragraph{Viability and set invariance.} The first-order safe-action set is the contingent-cone condition of viability theory \cite{Aubin1991,Aubin2011}, whose exactness under constraint qualifications goes back to Nagumo \cite{Nagumo1942}; the controlled-invariance literature for linear systems with polyhedral constraints \cite{Blanchini1999} computes the same objects by linear programming. Partial observation enters viability theory through guaranteed kernels and set-valued observers \cite{Cardaliaguet2007}, which give sufficient conditions; the common-action obstruction is the necessary first-order condition they must respect, and Theorem~\ref{thm:dual} is its certificate. Barrier certificates \cite{PrajnaJadbabaie2004} and control barrier functions \cite{Ames2017} certify safety; the measure certificate certifies the opposite, that no control satisfies the barrier condition on an atom, and it is sparse. Hamilton--Jacobi reachability \cite{Mitchell2005} computes the full-information value; the gap between $V^+$ and $V^-$ in Remark~\ref{rem:isaacs} is what partial information removes.

\paragraph{Minimax, Farkas and semi-infinite duality.} The use of Sion's theorem to turn an existence question into a measure is classical in semi-infinite programming \cite{HettichKortanek1993,AndersonNash1987} and in robust optimisation \cite{BenTalNemirovski1999,BenTal2009}; the specialisation to affine rows and a compact convex action set gives the exactness, and Helly's theorem gives the support bound that Carath\'eodory alone would miss by one. The scenario approach \cite{CalafioreCampi2006,CampiGaratti2008} finds the same number $k+1$ as the number of support constraints of a random convex programme; the connection is that both count the constraints a Helly argument needs. The adjustability question of Wei and Zhang \cite{WeiZhang2024}---when do adaptive decisions improve on static ones in a two-stage robust linear programme---is the open-loop gap of Table~\ref{tab:sweep} for linear data.

\paragraph{Dynamic games and information design.} The envelope of Section~\ref{sec:dynamic} is the finite partially observed dynamic programme \cite{Astrom1965,BertsekasShreve1978} written on review trees; upper and lower values are those of Isaacs \cite{Isaacs1965} and Ba\c{s}ar and Olsder \cite{BasarOlsder1999}. The robust value over couplings is a problem of information design \cite{BergemannMorris2019} with the coupling as the designer's instrument; the observation that the worst coupling can be interior to the class (Example~\ref{ex:sweep}) is the dynamic analogue of the mixing in Theorem~\ref{thm:dual}. The companion papers supply the continuous bridge \cite{Abaee2026cert}, the probabilistic layer \cite{Abaee2026prob} and the set-valued calculus \cite{Abaee2026calc} that the present dual serves.

\paragraph{What the certificate is for.} The reader who needs only the verdict can solve the primal linear programme. The certificate is for three other uses: it is \emph{portable} (a few rows and one support-function evaluation, checkable by a proof assistant); it is \emph{diagnostic} (support and facet multipliers run the three tests of Section~\ref{sec:taxonomy}); and it is \emph{prescriptive} with prices (the benchmark's five units of stock per unit of demand floor, the three-branch instance's $\tfr7{50}$). The same three uses carry over to the dynamic layer: the pair (coupling, recursion) is a portable dynamic certificate, the sweep is a diagnostic, and the open-loop gap prices the review.

\paragraph{Open problems.}
\begin{openproblem}[continuous review envelopes]\label{op:continuous}
Formulate the envelope of Definition~\ref{def:tree} for continuous state and observation spaces with action-dependent observations and prove the tree-supremum identity under measurable selection, including the case where the controller's actions determine which cells are observed; the finite theory here should be its exact restriction.
\end{openproblem}
\begin{openproblem}[robust envelopes beyond polytopes]\label{op:robust}
Extend Proposition~\ref{prop:robust} to classes of couplings defined by moment constraints or Wasserstein balls, and characterise when the worst coupling is extreme rather than interior.
\end{openproblem}
\begin{openproblem}[exact integrated duality]\label{op:integrated}
Identify conditions under which the inequality \eqref{eq:gammah} is tight for every refining schedule without passing through a mesh programme, i.e. a Sion-type theorem on the space of adapted controls with action-independent schedules; the three-branch instance shows the bound is tight at the threshold.
\end{openproblem}
\begin{openproblem}[non-convex residuals]\label{op:nonconvex}
For finite action sets, find a certificate notion for the residual obstruction (obstructed on $U$, not on $\conv U$) that is sparser than the full combinatorial check and sound; the parity instance is the smallest test.
\end{openproblem}
